% Notes of the Montreal seminar, 1970 (cote 129, typed part): archive cote 129, site page 70 (PDF page 71).
% typescript. Draft compared: none.
% AI-transcribed by Claude Opus 5.5 (claude-opus-5-5, adaptive thinking, effort high), Message Batches API, batch mtl129_full (msgbatch_01AGnDKHwy6ATCM7ehUa6263), 2026-11-13; not yet checked by eye.
\scanpage{129/70}
\mspage{\hand{III-15}}

on obtient une équivalence au lieu d'une anti-équivalence. Cette équi-%
valence peut s'expliciter comme suit:
\[
G \mapsto \underline{\mathrm{Hom}}_{gr}(\hand{\mu},G) \overset{\mathrm{def}}{=\!=} (\underline{\mathrm{Hom}}_{gr}(\hand{\mu}(n),G(n)))_{n\geq 1}
\]

6.6. On dit qu'un groupe de $BT$ $G$ sur $S$ est \uline{à fibres connexes}
si $G(1)$ est infinitésimal. Il revient au même de dire que tous les
$G(n)$ sont infinitésimaux.

Si toutes les caractéristiques résiduelles de $S$ sont égales
à $p$, les groupes de $BT$ toroidaux sont à fibres connexes (II,2.2).

On montrera plus loin que, si $p$ est localement nilpotent sur
$S$, un groupe de $BT$ est à fibres connexes si et seulement si il est
aussi un groupe de Lie formel.

Nous allons donner ici une première définition élémentaire des
groupes de Lie formels et étudier à quelles conditions un groupe de Lie
formel est un groupe de $BT$.

6.7. On appelle \uline{variété de Lie formelle ponctuée} (ou \uline{augmentée}) sur
$S$ un faisceau $X$ sur $S$ (pour la topologie $\hand{\mathcal{T}}$, (f.p.p.f.) $\leq \hand{\mathcal{T}} \leq$ (f.p.q.c.))
muni d'une section $S \overset{\hand{\xi}}{\to} X$ et tel que localement (pour la topologie de
Zariski) sur $X$, on ait
\[
X \simeq \varinjlim_{n} \operatorname{Spec} \underline{\hand{\mathcal{O}}}_S[[T_1,\ldots,T_r]]/(T_1,\ldots,T_r)^n .
\]

\uline{Remarques}. 1) On peut expliciter ce dernier isomorphisme en disant
que pour tout $S' \to S$, $S'$ quasi-compact, on a

\mspage{\hand{70}}
